% Chapter 34: Mirrors, Lenses, and Eyes
\chapter{Mirrors, Lenses, and Eyes}

\begin{mainline}

\step{A Counterintuitive Opening}

Start with something you have done ten thousand times but may not explain clearly: stand before a mirror and raise your left hand.

The person in the mirror raises a hand. Which one? Almost everyone immediately says ``the right'': the mirror reverses left and right. But think again: why does it not also reverse up and down? A flat mirror has no distinction between horizontal and vertical. Why would it single out ``left and right''? If it truly favors a direction, should laying the mirror horizontally on the floor make it reverse up and down instead? Try at home, but the result will disappoint you: however the mirror is placed, what gets ``reversed'' is always left and right, never up and down.

A second observation: many barbershop mirrors are only half a person's height and hang fairly low, yet stepping back a pace or two lets you see your entire body, head to soles. Intuition says a half-height mirror should fit at most half of you. The opposite is true, and here is something that will repeatedly trip you up: retreat another ten steps and the required mirror height does not increase by a single millimeter.

Another everyday puzzle: supermarkets hang bulging convex mirrors at aisle corners, and a car's passenger-side mirror carries a warning: ``Objects in mirror are closer than they appear.'' Why deliberately install a ``lying'' mirror? Is it necessarily lying?

The same physics underlies all three: light travels straight and changes direction at an interface, while our eyes and brains infer positions only by assuming that light arrived along straight lines. We will unpack that physics and explain your eyes, spectacles, cameras, microscopes, and telescopes, along with something more important: where this whole ``ray'' theory collapses.

Incidentally, the same backward inference occurs at water surfaces. On a windless lake, distant mountains and trees have perfectly faithful reflections. The water is a horizontal plane mirror, with virtual images suspended below it at depths equal to the real objects' heights. Wind breaks that surface into countless tiny mirrors tilted in different directions, fragmenting the reflection into dancing bands of light. From beginning to end, a reflection's ``clarity'' simply reports how smooth the surface is.

\step{Make a Guess First}

Put your intuition on the table before reading further; no secretly changing your bets afterward:

\begin{itemize}
  \item Guess one: a mirror image reverses left and right because the mirror does something horizontally that it does not do vertically. Agree? If so, what does it do?
  \item Guess two: how tall must a vertical plane mirror be for a 1.7-meter-tall person to see their whole body? How does this relate to body height, and does viewing distance matter?
  \item Guess three: point a convex lens, a magnifying glass, at a building outside. What will appear on white paper behind the lens: an upright or inverted image? Cover the lens's left half with your hand. Does the image's left half disappear, does the whole image dim, or does something else happen?
  \item Guess four: a nearsighted person cannot see distant objects clearly without spectacles. Does the eye converge light ``too strongly'' or ``not strongly enough''?
\end{itemize}

Write down your answers. Return at the chapter's end to mark your own paper.

\step{Hands-on Experiment}

\begin{experiment}[Half a Mirror, a Spoon, and a Sheet of Paper]

\textbf{Experiment one: the half-height mirror.} Find a full-length mirror and stand with your back to a window. Have a partner place two sticky notes on it: one just blocking the reflected light from the top of your head, the upper edge of your reflected head, and one blocking your soles. Measure the distance between the notes and your height. Walk two steps closer, then three steps back, and have your partner place them again. Their positions barely move, and their separation remains about half your height. Even retreating across the room changes nothing. This is the plane mirror's first piece of intelligence for you.

\textbf{Experiment two: a spoon.} Take a stainless-steel spoon. Look first at its bulging back: your face is upright and reduced at every distance. Turn it over and look into the hollow side: very close up, your face is upright and enlarged. Move the spoon slowly away; at one point the image becomes a blur, then suddenly flips upside down, shrinking as the distance increases. One curved surface, two behaviors, separated by a ``flipping distance.'' That distance is this chapter's central new term: focal length.

% BEGIN TRANSLATOR SAFETY NOTE
\textbf{Translator safety note:} Do not point the lens at the Sun or look toward it through optical devices. Keep this activity aimed at ordinary scenery, with the Sun excluded. See \href{https://science.nasa.gov/eclipses/safety/}{NASA solar-viewing safety guidance}.
% END TRANSLATOR SAFETY NOTE

\textbf{Experiment three: scenery on paper.} During daylight, close the curtains except for a narrow gap, letting bright outdoor scenery, buildings, trees, or streetlights, illuminate a corner of the room. Point a convex lens, reading glasses or a magnifying glass, outside. Hold white paper upright a few tens of centimeters behind it and move the paper back and forth. At one position, an outdoor scene appears: sharp, colored, and inverted both vertically and horizontally. This ``painting'' uses no pigment. Remove the paper and it remains in the air; put your hand there and spots of light fall on it. Rays really converge there. This is a real image, the thing you are meant to catch physically in this experiment. Complete guess three's test by covering half the lens. Half the image does not disappear; the whole image merely darkens. Record that counterintuitive result for explanation later.

\end{experiment}

\step{The Old Model Runs into Trouble}

Humanity once held an exceptionally long-lived model of seeing: the eyes extended a ``line of sight'' that touched objects like a hand, making them visible. Many ancient Greek thinkers genuinely believed this. Its simple advantage was to explain the active quality of ``I see wherever I look.'' Its problems were equally simple: in a pitch-dark room, perfectly healthy eyes cannot touch anything; uncover your eyes and seeing takes no perceptible time, so that ``hand'' must extend extraordinarily quickly. Mirrors are more fatal still: the reflected chair plainly is not behind the mirror, yet the ``sight-hand'' touches it there. When did it learn to pass through walls?

Replace this with Chapter 33's model: an object emits or reflects light in every direction, and you see it when a small bundle enters your pupil. Darkness is immediately explained, but an unfinished question remains: \emph{is an image a picture that lives somewhere?} Many people unconsciously imagine a mirror image as a ``picture pasted onto the surface.'' That model collapses just as quickly. Move left and the image's position shifts. Meanwhile, a friend on the other side looking into the same mirror sees the image at the same depth in space as you do. Two people viewing ``one picture stuck on the mirror'' from different directions could not agree on that depth. The image is not on the surface. It does not even ``sit'' anywhere waiting to be seen; it is an arrangement of rays established before they reach your eyes.

A deeper difficulty can wait here for later: if light travels strictly straight, each point in the real image should be determined by one ray. Yet experiment three showed that covering half the lens merely dims the image, so every image point is built by \emph{many} rays together. The ``one ray, one point'' picture must also be upgraded.

\step{Inventing New Concepts}

Like nineteenth-century instrument makers, let us construct the necessary concepts one at a time.

\textbf{First: object points and ray bundles.} Divide an object into many points, each emitting light in every direction. What enters your eye, reaches a mirror, or passes through a lens is always just a narrow cone of light from that point. To track ``seeing,'' follow how an optical device reshapes this cone.

\textbf{Second: real and virtual images.} If a cone from an object point reconverges after passing a mirror or lens, its rays actually intersect at a \emph{real image point}. The scenery on paper in experiment three is a collection of such points. If the cone becomes more divergent, but backward extensions of its rays meet at one point, that is a \emph{virtual image point}. Light never visited it, but your eyes cannot distinguish ``light diverging from there'' from ``light manipulated to seem to diverge from there.'' Eyes and brain recognize only the shape of the cone entering the pupil, so you ``see'' something that is not there. Your mirror reflection is a virtual image: not ``false'' in the sense of a hallucination, but ``absent'' in the sense of actual ray paths.

\textbf{Third: focus and focal length.} A bundle parallel to the principal axis strikes a concave mirror or convex lens and converges through, or appears to pass through, one point on that axis, the \emph{focus}. Its distance from the device is the \emph{focal length}. Focal length is an instrument's identity card, recording how strongly it bends light. For a spherical mirror and paraxial rays, it equals half the radius of curvature. For a lens, it depends on both surfaces' curvature and the glass's refractive index. The spoon's ``image-flipping distance'' is its concave face's focal length: when the object lies at the focus, reflected rays become parallel and the image retreats to infinity, leaving a blur.

\textbf{Fourth: three principal rays.} Predicting image position does not require tracking every ray in the cone, only three convenient ones: a ray parallel to the axis passes through the focus; one through the focus emerges parallel to the axis; and one through the lens center, or the spherical mirror's vertex, barely changes direction. Two locate the intersection, and the third checks it. The following two diagrams are the chapter's main working drawings; read them alongside the text.

With focus and focal length established, the spherical-mirror family's everyday jobs become clear. Concave mirrors converge: solar cookers and satellite dishes focus parallel signals collected over a large area onto a little receiver at the focus. Flashlights and headlights work backward, turning a source at the focus into a parallel beam. Makeup mirrors exploit the upright enlarged virtual image for $s<f$, enlarging your face at the cost of requiring it to remain within the focal length. Convex mirrors diverge, trading smaller images for a wider field of view at supermarket corners, garage exits, and cars' passenger sides. Two easily confused neighbors also deserve mention: a door peephole is a concave lens providing a wide-field, upright reduced virtual image; a dentist's little mirror is merely a plane mirror that redirects the view without magnifying it.

\begin{center}
\begin{tikzpicture}[scale=1.0,>=Stealth]
  \draw[thick] (0,-2.1)--(0,2.1);
  \fill[pattern=north east lines] (0,-2.1) rectangle (0.18,2.1);
  \node[right] at (0.1,-2.35) {Mirror};
  \draw[->,very thick] (-3,0)--(-3,1.5);
  \node[left] at (-3,0.75) {Object};
  \fill (-3,1.5) circle (1.2pt) node[above left]{$P$};
  \draw[->,very thick,dashed] (3,0)--(3,1.5);
  \node[right] at (3,0.75) {Virtual image};
  \fill (3,1.5) circle (1.2pt) node[above right]{$P'$};
  \draw[->] (-3,1.5)--(0,0.9);
  \draw[->] (0,0.9)--(-1.5,-0.6);
  \draw[->] (-3,1.5)--(0,0.1);
  \draw[->] (0,0.1)--(-1.5,-0.6);
  \draw[dashed] (0,0.9)--(3,1.5);
  \draw[dashed] (0,0.1)--(3,1.5);
  \draw (-1.5,-0.6) circle (0.14);
  \node[below] at (-1.5,-0.85) {Eye};
  \draw[dashed,gray] (-3,0)--(3,0);
  \node[below] at (0,-3.0) {\shortstack{Fig. 1: Plane mirror. Backward ray extensions meet at $P'$;\\object and image are mirror-symmetric.}};
\end{tikzpicture}
\end{center}

\begin{center}
\begin{tikzpicture}[scale=0.78,>=Stealth]
  \draw[->] (-7.2,0)--(7.2,0);
  \draw[thick] (0,-2.3)--(0,2.6);
  \draw[->,thick] (0,2.6)--(0.18,2.42);
  \draw[->,thick] (0,-2.3)--(-0.18,-2.12);
  \draw[->,thick] (0,2.6)--(-0.18,2.42);
  \draw[->,thick] (0,-2.3)--(0.18,-2.12);
  \fill (2,0) circle (1pt) node[below]{$F$};
  \fill (4,0) circle (1pt) node[below]{$2F$};
  \fill (-2,0) circle (1pt) node[below]{$F$};
  \fill (-4,0) circle (1pt) node[below]{$2F$};
  \draw[->,very thick] (-5,0)--(-5,1.6);
  \node[left] at (-5,0.8) {Object};
  \draw[->] (-5,1.6)--(0,1.6);
  \draw[->] (0,1.6)--(2,0);
  \draw (2,0)--(3.33,-1.07);
  \draw[->] (-5,1.6)--(0,0.53);
  \draw (0,0.53)--(3.33,-1.07);
  \draw[->,very thick] (3.33,0)--(3.33,-1.07);
  \node[right] at (3.33,-0.55) {Real image};
  \fill (3.33,-1.07) circle (1.2pt);
  \node[below] at (0,-2.9) {\shortstack{Fig. 2: Convex lens. Object beyond twice the focal length;\\image between $F$ and $2F$, inverted and reduced.}};
\end{tikzpicture}
\end{center}

\textbf{Fifth: what kind of machine is an eyeball?} Now unpack your own eye. Cornea and crystalline lens form a lens system that images the outside world onto the retina, its light-sensitive surface. The ``eye as camera'' analogy helps, but hear its qualifications first. \emph{Where it fits}: converging lenses in front, a detector behind, and a fixed lens-to-detector distance. Focusing changes the crystalline lens's curvature and focal length, analogous to extending a camera lens. Both form inverted, reduced real images. \emph{Where it fails}: the retina has no uniform pixel grid. A small central region, the macula, is packed with light-sensitive cells, while the sparse periphery mainly detects movement. Nor does the eye sit still and ``take a photograph''; it constantly jumps and scans. Much of your apparently stable high-resolution world is the brain's reconstruction from fragments. Using the camera analogy teaches optics; believing it literally teaches bad neuroscience.

A famous question immediately arises: if the retinal image is inverted, why does the world look upright? Once this chapter's concepts are clear, the question dissolves. There is no little observer inside who must inspect the retinal picture. Each receptor simply converts incident intensity and color into electrical signals. The brain receives a collection of numbered signal wires, not a photograph needing to be ``turned right-side up.'' ``Up'' and ``down'' are labels it assigns using gravity, posture, and feedback from the hands, with labeling rules entirely calibrated through experience. Psychologists have had people wear left-right-reversing or inverting spectacles for days. Initially the world is overturned; after several days the brain relabels its inputs and people cycle and pour water normally. Removing the spectacles requires adaptation again. Where optics forms an image, and whether it is inverted, is physics; your sense that the world is ``upright'' is labeling. Each handles its own part of the account.

\end{mainline}

\begin{mainline}

\step{The Model in One Sentence}

\begin{oneline}
Mirrors and lenses never ``manufacture'' images: they redirect each ray according to reflection and refraction. If redirected rays from one object point meet again, they form a real image; if only their backward extensions meet, they form a virtual image. Your eyes recognize light cones, not their place of origin.
\end{oneline}

\step{The Mathematical Translator}

Intuition is ready; give it a calculator. All imaging problems reduce to one equation and one magnification relation. First set the sign convention: object distance $s$ is positive for an object on the incident-light side; image distance $s'$ is positive for a real image, on the side actually reached by light, and negative for a virtual image. A converging device, convex lens or concave mirror, has positive focal length $f$; a diverging device, concave lens or convex mirror, has negative focal length. With this convention, thin lenses and spherical mirrors share one equation:

\[
\frac{1}{s}+\frac{1}{s'}=\frac{1}{f},
\qquad
m=-\frac{s'}{s}.
\]

The question is ``with an object at $s$, where and how large is the image?'' The sole device input $f$ describes its converging power. The image-to-object height ratio is $m$, with a negative sign indicating inversion.

As usual, subject the formula to extreme-case checkups before use:

\begin{itemize}
  \item \textbf{$s\to\infty$}: $1/s\to0$ gives $s'=f$. Parallel light converges at the focus, exactly the definition of focal length. The equation returns the definition without declaring independence. Burning a spot by concentrating sunlight with a magnifying glass uses this: the Sun is at infinity and its spot lies in the focal plane.

% BEGIN TRANSLATOR SAFETY NOTE
\textbf{Translator safety note:} Do not reproduce the burning example. Concentrated sunlight can injure eyes and ignite materials; keep lenses away from direct sunlight when not in controlled use. See \href{https://science.nasa.gov/eclipses/safety/}{NASA eye-safety guidance} and \href{https://www.london-fire.gov.uk/incidents/2022/january/flat-fire-bow/}{London Fire Brigade's sunlight-fire warning}.
% END TRANSLATOR SAFETY NOTE
  \item \textbf{$s=2f$}: $1/s'=1/f-1/(2f)=1/(2f)$, $s'=2f$, and $m=-1$. Equal size, inverted, with object and image symmetrically placed. Move Figure 2's object farther away and the image moves toward $F$, explaining why a camera photographing distant scenery focuses almost at the focal plane.
  \item \textbf{$s\to f$}: $1/s'\to0$ and $s'\to\infty$. An object at the focus sends the image to infinity: emerging rays are parallel. Reverse the arrangement and put a bulb at the focus to obtain a searchlight's or flashlight's parabolic reflector.
  \item \textbf{$s<f$}: $1/s'=1/f-1/s<0$, so $s'$ is negative: a virtual image, with $|m|>1$, upright and enlarged. This is a magnifying glass: light from each printed letter becomes more divergent after the lens, and your eyes infer a larger, more distant virtual letter.
  \item \textbf{$f<0$ (concave lens or convex mirror)}: for every positive $s$, $s'$ remains negative and $0<m<1$, always an upright, reduced virtual image. Convex mirrors compress a wide field into a small surface at the cost of smaller, seemingly more distant images. Hence the passenger-side mirror warning. It is not ``lying'' but exchanging size for field of view; your distance judgment needs recalibration.
  \item \textbf{$f\to\infty$}: $1/f\to0$, $s'=-s$, and $m=+1$. A ``lens'' with no converging power is a plane mirror: equal object and image distances, equal size, upright, virtual. The plane mirror is neatly incorporated as a limit of the same equation, a mark of good theory: old results do not need separate establishments.
\end{itemize}

Now calculate with real numbers for the eyes you are using. The lens-to-retina distance is about $1.7$ centimeters, a fixed $s'$. Looking far away, $s\to\infty$, the crystalline lens relaxes and $f=1.7$ centimeters. Reading at $25$ centimeters gives $1/f=1/0.25+1/0.017$, so $f\approx1.6$ centimeters. The ciliary muscle makes the lens more convex, shortening focal length by about $1$ millimeter to refocus from distant mountains to a book. This astonishingly small adjustment is also where trouble appears when the system cannot cope. A nearsighted eye is too long front to back, placing distant images \emph{in front of} the retina: guess four's answer is too much convergence, not too little. A concave spectacle lens first diverges the light slightly. Lens power is measured in diopters, $P=1/f$, with $f$ in meters. A $-2.00$-diopter lens has focal length $-0.5$ meters and forms a virtual image of an infinitely distant object $0.5$ meters before the eyes, within the nearsighted person's clear range. Farsightedness and presbyopia are the reverse: nearby images fall \emph{behind} the retina, requiring a convex lens to add convergence.

\textbf{Combining devices: cameras, microscopes, and telescopes.} The single-lens equation already lets us see through all three major optical instruments.

The camera is the most straightforward, directly applying Figure 2. For distant subjects, $s'\approx f$, and the sensor sits near the focal plane. For nearby subjects, $s$ decreases and $s'$ must increase, so the lens extends: that is the focusing motor's job. A variable aperture inside is described by its $f$-number, focal length divided by aperture diameter. Open it wider and more light enters in a broader cone, blurring out-of-focus points into larger spots and separating a sharp subject from a soft background. Stop down and foreground and background both fit within an acceptable range of sharpness. Stop down too far, for example below $f/16$, and each point's diffraction spot becomes larger than a pixel, softening the entire picture. Choosing aperture negotiates among light collection, depth of field, and diffraction blur, invoking all three of this chapter's boundaries.

A microscope is a two-stage relay. Its very short-focal-length objective, on the millimeter scale, forms an enlarged inverted real image of a specimen just outside its focus. The eyepiece then treats this image as its object, following the magnifying-glass rule, $s<f$ with an upright enlarged virtual image, and sends it to the eye. Total magnification is the product of the two stages. This explains the short, broad objective: its short focal length puts the specimen almost at the focus, while image distance extends to a dozen or more centimeters, making $m=-s'/s$ naturally large. Several hundredfold magnification sounds impressive, but each stage merely puts the same one-variable equation to work.

A telescope addresses a different problem: celestial objects are at infinity, so what it enlarges is not ``size'' but ``angular extent.'' A long-focal-length objective forms a real image in its focal plane; a short-focal-length eyepiece views it as a magnifier. Angular magnification is $M=f_{\text{obj}}/f_{\text{eye}}$. Changing to a shorter eyepiece changes magnification, which is why amateur astronomers keep rows of eyepieces. But remember resolution: magnification only enlarges the same spot; aperture $D$ determines whether fine detail can be resolved. Magnification beyond what the aperture supports is ``empty magnification'': the image grows, but detail does not. Cheap telescopes advertised at ``six hundred times'' mostly offer empty magnification. Experts ask aperture first, not magnification.

\end{mainline}

\begin{mathbridge}
\textbf{Where the thin-lens equation comes from.} Only Figure 2 and two pairs of similar triangles are needed. Let object height be $h$ and image height $h'$. First pair: the ray through the optical center is undeviated, so object tip, lens center, and image tip are collinear, giving
\[
\frac{h'}{h}=\frac{s'}{s}.
\]
Second pair: the ray incident parallel to the axis passes through the focus. Follow it: at the lens plane its height is $h$, then it heads directly toward focus $F$ and continues to the image plane at height $h'$. Similar triangles on either side of the focus give
\[
\frac{h}{h'}=\frac{f}{s'-f}.
\]
Multiply the equations to eliminate $h'/h$:
\[
1=\frac{s'}{s}\cdot\frac{f}{s'-f}
\;\Longrightarrow\;
s(s'-f)=s'f
\;\Longrightarrow\;
ss'-sf=s'f.
\]
Dividing by $ss'f$ gives exactly $\dfrac{1}{f}-\dfrac{1}{s'}=\dfrac{1}{s}$, which rearranges to the imaging equation. Two assumptions slipped in: the lens is thin enough to treat all refraction as occurring in one plane, and rays lie close enough to the axis to use ``height'' directly as ``arc length.'' The next section examines how trustworthy those assumptions are.

\textbf{Why diopters add.} Two thin lenses in contact have total power $P=P_1+P_2$. The first changes the cone's convergence or divergence by $P_1$, in inverse meters, and the second adds another change $P_2$ to the outgoing shape. This licenses the optometrist's practice of stacking trial lenses: it is powers that add, not focal lengths.
\end{mathbridge}

\begin{mainline}

\step{The Model's Limits}

Every formula must state its jurisdiction. Ray-based imaging has three boundaries.

\textbf{First: thin lenses and paraxial rays.} The derivation assumed rays close to the principal axis and traveling at small angles. A wide real lens bends marginal rays strongly, breaking the paraxial approximation: rays near a spherical lens's edge and center do not converge at the same point, producing \emph{spherical aberration}. History's costliest example was the Hubble Space Telescope. Its 2.4-meter primary mirror differed from the intended shape near its edge by only about 2 micrometers, less than a thirtieth of a hair's diameter, yet this spread starlight into blurry spots. Sharpness returned only when astronauts fitted ``spectacles,'' corrective optics, in 1993. Anyone tempted to think ``close enough is good enough'' should read three times the story of 2 micrometers ruining the sharpness of a billion-dollar instrument.

\textbf{Second: color.} Glass has slightly different refractive indices for different colors, so a lens has different focal lengths for blue and red, producing colored image fringes: \emph{chromatic aberration}. Frustrated by it, Newton bypassed refraction and designed a reflecting telescope with a concave mirror, which treats all colors equally. All today's largest optical telescopes are reflectors, his legacy. Spherical and chromatic aberration are merely the family's best-known members. Oblique bundles can grow comet-like tails, coma; horizontal and vertical focal lines can fail to coincide, astigmatism; and straight lines can curve into barrel or pincushion shapes, distortion. The seven, eight, or dozen-plus elements in a camera lens mostly cancel one another's aberrations rather than ``magnify.'' One element's errors are paid back by another with an opposing shape.

\textbf{Third and deepest: light does not always travel straight.} Return to experiment three: covering half the lens leaves the image complete but dimmer. Each image point is built by a whole cone from its object point, so half the light merely halves every point's brightness. That is understandable, but also shows that ``light from every point independently travels straight'' is not a secure picture. Tiny holes bring the wall down. Pinhole-camera makers know that smaller holes geometrically sharpen images, but below a few tenths of a millimeter, images blur again: emerging light spreads instead of staying straight. This is neither a lens defect nor poor workmanship. The axiom of straight propagation itself fails when aperture approaches light's wavelength, about $0.5$ micrometers for visible light. Chapter 35 rebuilds the picture with waves. For now, note that pinhole imaging has an optimal aperture: geometrical blur grows with diameter, diffraction blur grows as diameter shrinks. A ten-centimeter-deep camera works best around $0.2$ to $0.3$ millimeters, a real-world compromise between competing trends.

This boundary also has an everyday consequence: \emph{resolution is limited}. Through any circular aperture of diameter $D$, pupil, lens, or telescope primary, a point source forms a finite spot. Two sources separated by less than about $\theta\sim\lambda/D$ cannot be distinguished. Human pupils are roughly $3$ to $5$ millimeters across. Taking visible wavelength $5.5\times10^{-7}$ meters gives $\theta\sim10^{-4}$ radians, about half an arcminute. Normal visual acuity of ``$1.0$'' on an eye chart corresponds to resolving $1$ arcminute, and receptor spacing in the macula matches that density. The pupil's diffraction limit and retinal sampling density meet at the same order of magnitude; evolution wasted neither. This is not praise of coincidence but a clue: your eye's optical design already reaches the physical ceiling. Finer detail requires a larger aperture, a telescope, or a shorter wavelength.

\end{mainline}

\begin{challenge}
\textbf{The Rayleigh criterion and its 1.22.} More precisely, the resolution limit is $\theta_{\min}\approx1.22\lambda/D$. The coefficient $1.22$ comes from the first dark ring in a circular aperture's diffraction pattern, described by a Bessel function with its first zero at $x\approx3.83$, giving $3.83/\pi\approx1.22$. Chapter 35 supplies the wave-optics derivation. Here the logical point matters: $1.22$ is not an empirical fit, but the definite answer to ``a wave passing through a circular hole.'' For microscopy, the corresponding lateral resolution is $d\sim\lambda/(2\,\mathrm{NA})$, where numerical aperture $\mathrm{NA}=n\sin\theta$ describes the breadth of the collected cone. Oil-immersion objectives raise $n$ from $1$ to about $1.5$, exploiting the numerator. Visible-light microscopes stop at roughly $0.2$ micrometers. Smaller things are not hidden by inadequate magnification, but because visible light itself is too coarse a probe. Electron microscopes can see viruses and crystal lattices because accelerated electrons have de Broglie wavelengths down to picometers. Optical microscopy's ceiling is the first knock on the door of the quantum section's ``matter is also a wave.''

\textbf{Reflecting telescopes' aperture account.} Reverse $\theta\sim\lambda/D$: resolving two celestial objects separated by $\theta$ requires aperture $D$ at least $\lambda/\theta$. That is why radio telescopes span hundreds of meters, not out of greater ambition for magnification, but to pay the aperture tax imposed by centimeter wavelengths for the same resolution.
\end{challenge}

\begin{mainline}

\step{Explaining the Opening Phenomenon}

Now mark your answers, starting with stubborn ``left-right reversal.''

Study Figure 1. The mirror does just one thing: redirects every ray by the reflection law. For a vertical mirror, the line joining object point $P$ and image point $P'$ is \emph{perpendicular to the mirror}, with equal distances on either side. The reversed direction is neither left-right nor up-down, but \emph{front-back}, the direction into the mirror along your sightline. Your nose is 30 centimeters in front, its image 30 centimeters behind; your back is 30.4 centimeters in front, its image 30.4 centimeters behind. The front-back order is completely reversed. Where does apparent left-right reversal come from? Your comparison: you imagine going behind the mirror and \emph{turning around} to face your present self, then find that your raised left hand corresponds to that person's right. The crucial step is this imagined rotation about a vertical axis. Imagine instead a somersault about a horizontal axis, head down and feet up, and comparison with the reflected person reverses up and down while left and right agree. Your reflection is your opposite only front-to-back; whether you call it left-right or up-down reversal depends on how you imagine it turning around. The mirror treats horizontal and vertical equally. Your imagined turn favors the vertical axis.

Now the half-height mirror. Draw a line from your eyes to the virtual image's head: this is the direction of the actual reflected ray leaving the mirror. Draw another to its soles. Similar triangles show that the segment between their mirror intersections is exactly half your height, with no viewing distance appearing in the ratio. Retreat and the virtual image retreats proportionally, leaving the intersection spacing unchanged. That explains the stationary sticky notes. A mirror half your height and hung correctly always fits your full body.

The car-mirror warning is also settled. A convex mirror has $f<0$ and always forms an upright reduced virtual image. Estimating distance by image size, the brain places the following car too far away. Engineers exchange image size for a broad view, printing a warning to pay your intuition's tuition.

The spoon experiment is settled too. Its hollow side is a little concave mirror with a focal length of a few centimeters. Bring your face inside that distance, $s<f$, and see an upright enlarged virtual image. Move near the focus and reflected light approaches parallel, blurring the image. Farther away, $s>f$, rays genuinely converge into an inverted real image, suspended in the air between you and the spoon. Your eyes lie beyond it and receive it as an object at a distance. For the convex side, $f<0$; the formula keeps $s'$ negative and $m$ between $0$ and $1$, reliably upright and reduced. One utensil demonstrates the whole imaging equation.

One small bonus remains: the reflected you is your past self. If you stand 1 meter from a mirror, the round-trip light path is 2 meters, taking about $7$ nanoseconds. Looking into your reflection's eyes, you always see what the other looked like $7$ nanoseconds earlier. Absurdly tiny though that delay is, it belongs to a larger fact: no seeing is instantaneous. Every view is a receipt for a path traveled and time spent by light.

\step{Thought Experiments and Challenges}

\begin{thought}[Photograph the Flag on the Moon]

People sometimes claim a telescope can see the flag astronauts planted on the Moon. Weigh that claim with this chapter's limit. The unfolded flag is about 1 meter wide and the Moon about $3.8\times10^{8}$ meters away, subtending roughly $2.6\times10^{-9}$ radians. Resolving it requires aperture at least $D\sim\lambda/\theta\approx5.5\times10^{-7}/2.6\times10^{-9}\approx200$ meters, even optimistically omitting the factor $1.22$. Today's largest optical telescopes have apertures of order 10 meters, more than twenty times too small; Hubble does worse. So nobody has photographed that flag, nor will a single mirror do so in the future. But reverse the thought: the limit comes from wavelength divided by aperture, not from ``insufficiently good technology.'' Humanity has indeed seen lunar landing remains, photographed by Lunar Reconnaissance Orbiter from about 50 kilometers above the Moon. Reducing object distance from $3.8\times10^{8}$ to $5\times10^{4}$ meters immediately reduces the required aperture to centimeters. Engineers retried a case condemned by physics in a different court: if you cannot beat diffraction, move closer.

One more layer deserves inspection. Diffraction limits are evidence of light's wave behavior, yet they have already collected tax in a geometrical-optics chapter. Ray and wave models are not interchangeable worldviews to choose at whim: the former approximates the latter when aperture greatly exceeds wavelength. Every formula here lives within that approximation, and its failure modes, spreading through small holes and image points becoming spots, betray next chapter's truth.
\end{thought}

\begin{puzzles}
\item A periscope uses two parallel plane mirrors to turn light twice into the eyes. Does writing viewed through it appear normally or reversed? What about three reflections from three mirrors? (Hint: each plane reflection reverses front and back. Count reversals and consider what happens to handedness after an even number.)
\item Immerse a magnifying glass completely in water and inspect stones on the bottom. Is its magnification stronger or weaker than in air? (Hint: convergence depends on the refractive-index difference between glass and its surroundings. Water's 1.33 is closer to glass than air's 1.00. What happens to focal length? Then consider why a fish's crystalline lens is much more convex than a human's.)
\item In experiment three, covering half the convex lens dims the image without removing anything. Now replace the lens's \emph{upper half} with part of another lens of slightly longer focal length, leaving the lower half unchanged. What appears on the paper? (Hint: return to ``each image point is built by an entire cone.'' The cone has now split into two independently converging bundles.)
\end{puzzles}

End with the book's three recurring questions. What is the system's present state? A set of devices, mirrors, lenses, and eyeballs, plus light cones from object points, each reshaped into some converging or diverging form. What governs change? Reflection and refraction, unified in Chapter 33 by Fermat's stationary-time principle, plus an imaging equation compressing them to their simplest limit. How do measurements tell us? Put white paper at the image to test whether light genuinely converges, measure sticky-note positions on a mirror, and cover half a lens to test image completeness. Next chapter reveals that even the oldest axiom, ``light travels straight,'' is merely a sufficiently good approximation at certain scales. Behind every mirror in physics stands another mirror.

\end{mainline}
